\documentclass[11pt]{article}
\usepackage[spanish]{babel}
\usepackage[utf8]{inputenc}
\usepackage[T1]{fontenc}
\usepackage{geometry}
\geometry{margin=2.5cm}
\usepackage{enumitem}
\usepackage{natbib}
\usepackage{hyperref}
\hypersetup{colorlinks=true, linkcolor=blue, citecolor=blue, urlcolor=blue}

\title{Revisión de literatura\\
\large Dependencia del empleo y la economía peruana respecto al comercio\\
internacional de commodities: ¿hay un vacío que llenar?}
\author{}
\date{Julio 2026}

\begin{document}
\maketitle

\begin{abstract}
\noindent Esta revisión ordena la evidencia sobre cómo los shocks de precios internacionales
de commodities se transmiten a las economías locales y a los mercados de trabajo, con foco
en el Perú y sus regiones extractivas. Cubre cuatro bloques: (i) la transmisión de booms de
commodities a economías locales; (ii) la evidencia peruana sobre minería, ingreso, pobreza y
renta fiscal (canon); (iii) la informalidad como margen de ajuste laboral ante shocks
comerciales; y (iv) las herramientas de identificación (\textit{shift-share}/Bartik) que
hacen creíble este tipo de estimación. El objetivo es evaluar si la idea de investigación es
\emph{interesante}, \emph{factible} y si \emph{cubre un vacío}. Concluyo que el vacío más
claro y valioso no está en volver a medir el efecto de un boom sobre el ingreso —bien
documentado— sino en \textbf{la reasignación de la estructura del empleo (sector y
formalidad) ante la volatilidad de precios, distinguiendo auges de caídas (asimetría) y
canales de mercado laboral vs.\ renta fiscal}. Al final se propone una idea concreta.
\end{abstract}

\section{Motivación y pregunta que guía la revisión}
El Perú es una economía pequeña y abierta cuya canasta exportadora está fuertemente
concentrada en minerales (cobre, oro, zinc). Entre 2001 y 2013 las exportaciones mineras se
multiplicaron por un factor cercano a 7,5, lo que convierte al país en un laboratorio natural
para estudiar la dependencia del empleo y del ingreso respecto a los términos de intercambio.
La pregunta que organiza esta revisión es: \emph{¿qué sabemos ya, y qué no, sobre cómo los
shocks de precios internacionales de commodities se transmiten al mercado laboral de las
regiones extractivas del Perú?}

\section{Bloque 1. Transmisión de shocks de precios de commodities a economías locales}

La literatura moderna abandonó el enfoque puramente macro (la ``maldición de los recursos''
a nivel país, con resultados frágiles) y se movió hacia diseños \textbf{subnacionales} que
comparan localidades con distinta exposición a un mismo shock de precios internacional. Dos
trabajos son referencia obligada.

\citet{dube2013} (\textit{Review of Economic Studies}) explotan shocks exógenos de precios
internacionales sobre los principales bienes de exportación de Colombia (café y petróleo) y
muestran que el signo del efecto depende de la naturaleza del bien: una caída del precio del
café —intensivo en trabajo— reduce salarios y aumenta la violencia (efecto costo de
oportunidad), mientras que un alza del precio del petróleo —intensivo en capital— eleva
rentas y también la violencia (efecto ``rapiña''). La lección metodológica y conceptual es
central: \emph{no todos los commodities transmiten los shocks por el mismo canal}, y la
intensidad factorial del bien determina si el efecto pasa sobre todo por el mercado laboral o
por las rentas.

\citet{benguria2018} (NBER WP) estudian la transmisión de los ``super-ciclos'' de precios de
commodities en economías en desarrollo e incorporan explícitamente la \textbf{informalidad},
la movilidad interregional y mecanismos de economía política como parte del mecanismo de
transmisión. Trabajo posterior de los mismos autores \citep{benguria2024} documenta que los
efectos de un boom son \textbf{espacialmente desiguales} y dependen del nivel de calificación
del trabajador y de la región. Este es el marco teórico-empírico más cercano a la idea que se
quiere desarrollar.

\medskip
\noindent\textbf{Qué se llevan de este bloque:} (a) el diseño correcto es subnacional y de
exposición diferencial; (b) el canal (salarios vs.\ rentas) depende de la intensidad
factorial del commodity; (c) la informalidad y la movilidad son parte del mecanismo, no un
detalle.

\section{Bloque 2. Evidencia peruana: minería, ingreso, pobreza y renta fiscal}

\subsection{Encadenamientos productivos e ingreso local}
El punto de partida es \citet{aragon2013} (\textit{American Economic Journal: Economic
Policy}), que estudian la mina de oro Yanacocha con datos de hogares 1997--2006. Encuentran un
efecto positivo de la demanda de insumos locales de la mina sobre el ingreso real,
concentrado en el mercado proveedor y su entorno, alcanzando a trabajadores no calificados de
sectores \emph{no} mineros, junto con un aumento del precio de los bienes no transables. Su
mensaje —los encadenamientos hacia atrás de la industria extractiva pueden generar
externalidades positivas— fue muy influyente.

Sin embargo, \citet{aragon2020} (\textit{Environment and Development Economics}) revisitan el
caso y ofrecen una lectura más matizada: los efectos \textbf{se desvanecen con el agotamiento
de la mina} y responden a una dinámica de auge-y-caída del consumo. Esta tensión
(¿efecto permanente o transitorio?) sigue abierta y es directamente relevante para una
pregunta centrada en volatilidad de precios.

\subsection{Pobreza, desigualdad y el canon minero}
\citet{loayza2016} (\textit{World Development}) documentan que los distritos beneficiados por
el boom minero y las transferencias de canon redujeron la pobreza, pero con aumentos de
desigualdad y efectos heterogéneos según la cercanía a la producción. El canon abre un
\textbf{segundo canal} —el fiscal— distinto del mercado laboral: el Estado transfiere a los
gobiernos subnacionales una fracción de los impuestos mineros. La evidencia sobre su
eficacia es mixta: mientras algunos trabajos hallan mejoras en capital humano
\citep{aguero2017}, otros cuestionan la eficiencia y la sostenibilidad del gasto y sugieren
efectos de largo plazo cercanos a nulos \citep{aragon2023}. Para la idea de investigación,
esto implica que \emph{cualquier estimación del efecto de un shock de precios debe distinguir
el canal laboral privado del canal fiscal-canon}, o al menos discutir su confusión potencial.

\subsection{Calidad del empleo}
Un frente emergente —y aún poco poblado— examina no el \emph{nivel} sino la \emph{calidad}
del empleo generado por la minería en el Perú (condiciones, formalidad, estabilidad), usando
métodos multidimensionales \citep{jambrina2022}. Es precisamente aquí donde la evidencia es
más delgada.

\section{Bloque 3. La informalidad como margen de ajuste}

Fuera del Perú, la literatura de comercio y mercados laborales locales muestra que la
informalidad es un \textbf{margen de ajuste central} ante shocks comerciales.
\citet{dixcarneiro2017} (\textit{American Economic Review}), en el caso paradigmático de la
liberalización comercial brasileña, muestran que las regiones más expuestas sufren caídas
prolongadas del empleo formal y de los salarios, con una dinámica que se \emph{amplifica} en
el tiempo por la movilidad interregional imperfecta y el ajuste lento del capital. En trabajo
relacionado, los mismos autores documentan que los tránsitos hacia el desempleo y hacia la
\textbf{informalidad} son márgenes empíricamente grandes del ajuste al comercio.

La contraparte de precios de recursos existe pero es más escasa: por ejemplo, evidencia sobre
el \emph{bust} del precio del crudo en Ecuador muestra deterioro del mercado laboral local
\citep{idb_ecuador}. La asimetría auge/caída —¿la informalidad absorbe los malos tiempos
más de lo que los buenos la reducen?— está poco explorada para shocks de precios de
commodities y casi inexplorada para el Perú minero.

\medskip
\noindent\textbf{Qué se llevan de este bloque:} la informalidad no es solo un resultado de
interés, es el principal \emph{mecanismo de amortiguación} del mercado laboral. Estudiarla
como variable de estructura del empleo (y no solo el ingreso promedio) es donde la literatura
comparada ha encontrado los hallazgos más ricos.

\section{Bloque 4. Identificación: diseños \textit{shift-share} / Bartik}

La factibilidad econométrica descansa en el instrumento de participación-desplazamiento
(\textit{shift-share} o Bartik): la exposición regional se construye combinando la
composición productiva \emph{previa} de cada región (participaciones, \textit{shares}) con los
shocks comunes de precios internacionales (\textit{shifts}). La literatura reciente formalizó
sus supuestos de identificación desde dos ángulos complementarios:
\citet{goldsmith2020} (\textit{American Economic Review}) derivan la validez desde la
exogeneidad de las \emph{participaciones}, mientras que \citet{borusyak2022}
(\textit{Review of Economic Studies}) y \citet{adao2019} (\textit{Quarterly Journal of
Economics}) lo hacen desde la exogeneidad de los \emph{shocks} y corrigen la inferencia por la
correlación entre regiones con exposición similar. Para el caso peruano, las
\emph{participaciones} serían la especialización minera pre-shock de cada región/provincia
(p.\ ej.\ VAB minero, empleo minero o presencia de yacimientos en un año base) y los
\emph{shifts} los precios internacionales de cada metal. Esto hace la estrategia
\textbf{factible y estándar}, con datos públicos.

\section{Diagnóstico: ¿interesante, factible y con vacío?}

\begin{description}
  \item[¿Interesante?] Sí. La dependencia del empleo peruano respecto a los términos de
  intercambio es de primer orden para la política macro-laboral y fiscal, y conecta con
  debates vivos (volatilidad, informalidad, diversificación productiva).
  \item[¿Factible?] Sí. Existen (a) datos de hogares con panel rotativo y medición de
  informalidad (ENAHO), (b) precios internacionales de metales de alta frecuencia y (c) un
  diseño de identificación consolidado (Bartik). El principal reto es aislar el canal laboral
  del canal fiscal (canon).
  \item[¿Cubre un vacío?] \textbf{Parcialmente, y ahí está la oportunidad.} Lo que ya está
  bien establecido: el \emph{nivel} de ingreso/pobreza sube en el auge minero
  \citep{aragon2013,loayza2016}. Lo que está \emph{poco} establecido para el Perú:
  \begin{enumerate}[label=(\roman*)]
    \item la reasignación de la \textbf{estructura} del empleo (sector y formalidad), no solo
      el ingreso medio;
    \item la \textbf{asimetría} entre auges y caídas de precios (la informalidad como
      amortiguador de los \emph{busts});
    \item la separación limpia entre el \textbf{canal de mercado laboral} y el
      \textbf{canal fiscal (canon)};
    \item el papel de la coexistencia de \textbf{gran minería formal y minería
      informal/artesanal} en la respuesta del empleo.
  \end{enumerate}
\end{description}

\section{Idea de investigación sugerida}

\begin{quote}
\textbf{``Precios de metales y estructura del empleo en el Perú: reasignación sectorial e
informalidad como amortiguador ante la volatilidad de los términos de intercambio''.}

\medskip
\noindent\emph{Pregunta.} ¿Cómo modifican los shocks de precios internacionales de metales la
\textbf{composición sectorial y de formalidad} del empleo de los hogares en las regiones
peruanas, y en qué medida ese ajuste es \textbf{asimétrico} entre auges y caídas de precios?
\end{quote}

\noindent\textbf{Diseño propuesto (resumen).}
\begin{itemize}[leftmargin=1.4em]
  \item \emph{Shock:} índice de precios de metales específico por región, construido como
  \textit{shift-share} = (participación pre-shock de cada metal en la producción/empleo de la
  región) $\times$ (precio internacional del metal), siguiendo \citet{borusyak2022}.
  \item \emph{Outcome:} probabilidades de estar en cada celda (sector minero / otro transable
  / no transable) $\times$ (formal / informal), vía modelo multinomial o LPM, además de
  ingreso laboral como resultado secundario.
  \item \emph{Datos:} ENAHO (panel rotativo) 2004--2023; precios de metales (Banco Mundial /
  BCRP / LME); mapa de exposición minera pre-2004.
  \item \emph{Valor agregado:} (1) separar auges de caídas para testear asimetría; (2) incluir
  transferencias de canon como control/canal para aislar el efecto de mercado laboral del
  fiscal; (3) explotar la heterogeneidad gran minería vs.\ minería informal.
\end{itemize}

\noindent\textbf{Contribución esperada.} Pasar de ``el boom minero subió el ingreso'' a
``\emph{cómo} la volatilidad de precios reasigna a los trabajadores entre sectores y entre
formalidad e informalidad, y si ese ajuste protege o precariza a los hogares en las caídas''.
Esto llena el vacío señalado en el diagnóstico y dialoga directamente con
\citet{dixcarneiro2017} y \citet{benguria2024}, llevándolos al terreno —poco cubierto— de los
shocks de precios de recursos en una economía minera.

\bigskip
\noindent\textit{Nota sobre las referencias: los datos bibliográficos (revista, volumen, año)
fueron verificados en fuentes públicas durante la elaboración de esta revisión; conviene
confirmar números de página y ediciones finales antes de citar en un manuscrito formal.}

\begin{thebibliography}{99}
\setlength{\itemsep}{2pt}

\bibitem[Adão et al.(2019)]{adao2019}
Adão, R., Kolesár, M. y Morales, E. (2019). Shift-Share Designs: Theory and Inference.
\textit{Quarterly Journal of Economics}, 134(4), 1949--2010.

\bibitem[Agüero y Balcázar(2017)]{aguero2017}
Agüero, J. M. y Balcázar, C. F. (2017). The Value of Redistribution: Natural Resources and the
Formation of Human Capital under Weak Institutions. \textit{IZA Discussion Paper} No.\ 10884
(publicado luego en \textit{Journal of Development Economics}).

\bibitem[Aragón y Rud(2013)]{aragon2013}
Aragón, F. M. y Rud, J. P. (2013). Natural Resources and Local Communities: Evidence from a
Peruvian Gold Mine. \textit{American Economic Journal: Economic Policy}, 5(2), 1--25.

\bibitem[Aragón y Rud(2020)]{aragon2020}
Aragón, F. M. y Rud, J. P. (2020). Fading Local Effects: Boom and Bust Evidence from a
Peruvian Gold Mine. \textit{Environment and Development Economics}, 25(2), 182--203.

\bibitem[Aragón y Winkler(2023)]{aragon2023}
Aragón, F. M. y Winkler, H. (2023). The Long-Term Impact of a Resource-Based Fiscal Windfall:
Evidence from the Peruvian Canon. \textit{Resources Policy} (Elsevier).

\bibitem[Benguria et al.(2018)]{benguria2018}
Benguria, F., Saffie, F. y Urzúa, S. (2018). The Transmission of Commodity Price Super-Cycles.
\textit{NBER Working Paper} No.\ 24560.

\bibitem[Benguria et al.(2024)]{benguria2024}
Benguria, F., Saffie, F. y Zhang, S. (2024). Spatial Linkages and the Uneven Effects of a
Commodity Boom. \textit{Working Paper}.

\bibitem[Borusyak et al.(2022)]{borusyak2022}
Borusyak, K., Hull, P. y Jaravel, X. (2022). Quasi-Experimental Shift-Share Research Designs.
\textit{Review of Economic Studies}, 89(1), 181--213.

\bibitem[Dix-Carneiro y Kovak(2017)]{dixcarneiro2017}
Dix-Carneiro, R. y Kovak, B. K. (2017). Trade Liberalization and Regional Dynamics.
\textit{American Economic Review}, 107(10), 2908--2946.

\bibitem[Dube y Vargas(2013)]{dube2013}
Dube, O. y Vargas, J. F. (2013). Commodity Price Shocks and Civil Conflict: Evidence from
Colombia. \textit{Review of Economic Studies}, 80(4), 1384--1421.

\bibitem[Goldsmith-Pinkham et al.(2020)]{goldsmith2020}
Goldsmith-Pinkham, P., Sorkin, I. y Swift, H. (2020). Bartik Instruments: What, When, Why, and
How. \textit{American Economic Review}, 110(8), 2586--2624.

\bibitem[BID/Ecuador(2021)]{idb_ecuador}
Inter-American Development Bank (2021). The Labor Market Worsening Effects of a Resource Bust:
Evidence from the Crude Oil Price Shock in Ecuador. \textit{IDB Working Paper} Series.

\bibitem[Jambrina(2022)]{jambrina2022}
Jambrina-Canseco, B. (2022). Commodity Price Shocks and Quality of Employment in Peru.
\textit{Working Paper}, Oxford Poverty and Human Development Initiative (OPHI).

\bibitem[Loayza y Rigolini(2016)]{loayza2016}
Loayza, N. y Rigolini, J. (2016). The Local Impact of Mining on Poverty and Inequality:
Evidence from the Commodity Boom in Peru. \textit{World Development}, 84, 219--234.

\end{thebibliography}

\end{document}
